\documentclass[a4paper]{IEEEtran}
\usepackage[utf8]{inputenc}
\usepackage[T1]{fontenc}
\usepackage[english]{babel}
\usepackage{graphicx}
\usepackage{amsmath}
\usepackage{amssymb}
\usepackage{mathtools}
\usepackage{subcaption}
\usepackage[hidelinks]{hyperref}
\mathtoolsset{showonlyrefs}
\markboth{Seminar: Modern Methods of Information Processing, Winter Semester 2026/27}%
         {Seminar: Modern Methods of Information Processing, Winter Semester 2026/27}
\title{Diffusion models for trajectory planning}
\author{Nick~Bermbach}
\begin{document}
\maketitle

\section{Introduction}
\begin{itemize}
    \item Trajectory planning with multiple feasible routes, goals, and obstacle constraints
    \item Classical trajectory optimization versus learning a distribution of plausible trajectories
    \item Diffuser's architecture and conditioning, diffusion-based control policies, and a 2D planning study
\end{itemize}

\section{Problem formulation and diffusion basics}
\subsection{Trajectory planning}
\begin{itemize}
    \item State-action trajectories, system dynamics, and planning horizon; classical trajectory optimization as background~\cite{kelly2017trajectory}
    \item Start and goal conditions; collision avoidance, path length, and smoothness
    \item Classical optimization as a comparison: CMA-ES on a parameterized trajectory~\cite{hansen2016cma}
\end{itemize}
\subsection{Denoising diffusion models}
\begin{itemize}
    \item Forward noising, learned reverse process, and noise-prediction objective (DDPM)~\cite{ho2020ddpm}
    \item Distinction between physical trajectory time and diffusion timesteps
    \item Iterative sampling cost; DDIM as an optional accelerated sampler~\cite{song2021ddim}
\end{itemize}

\section{Diffuser: learning and conditioning trajectory distributions}
\subsection{Representation, architecture, and training}
\begin{itemize}
    \item Joint state-action trajectories; temporal U-Net and parallel denoising across the planning horizon~\cite{janner2022diffuser}
    \item Learning from offline trajectory data; multimodality and dependence on data coverage
\end{itemize}
\subsection{Goal conditioning and planning objectives}
\begin{itemize}
    \item Inpainting: fixing the current state and specified goals during iterative denoising
    \item Reward guidance using a separately learned return predictor; modifying the denoising distribution
    \item Planning as inference: $\widetilde p_\theta(\tau)\propto p_\theta(\tau)h(\tau)$~\cite{levine2018inference}
    \item Maze2D navigation with start-goal inpainting, in contrast to reward-guided locomotion
\end{itemize}
\subsection{Constraints and limitations}
\begin{itemize}
    \item Obstacle avoidance learned from trajectory data versus explicit test-time cost guidance
    \item Inpainting of specified values versus soft guidance of trajectory costs
    \item No general guarantees for collision avoidance or consistency with system dynamics
\end{itemize}

\section{From Diffuser to diffusion policies}
\subsection{Diffusion Policy}
\begin{itemize}
    \item Observation-conditioned generation of action sequences instead of joint state-action trajectories~\cite{chi2025diffusionpolicy}
    \item Behavior cloning from demonstrations; visual conditioning with temporal CNN or transformer
    \item Receding-horizon execution and the trade-off between action consistency and responsiveness
\end{itemize}
\subsection{One-Step Diffusion Policy}
\begin{itemize}
    \item Distillation of a pretrained diffusion-policy teacher into a single-pass action generator~\cite{wang2025onestep}
    \item Distribution matching using reverse-KL/score-difference objectives
    \item One-step inference, not one-step training; teacher and distillation costs remain
\end{itemize}
\subsection{Comparison of the three methods}
\begin{itemize}
    \item Generated variables, conditioning inputs, training data, and inference procedures
    \item Planning versus policy learning; use of receding-horizon execution in Diffuser and Diffusion Policy
    \item Flexibility, multimodality, runtime, and applicability to the 2D experiment
\end{itemize}

\section{Experimental study: 2D trajectory planning}
\subsection{Environment and training data}
\begin{itemize}
    \item Point robot with single-integrator dynamics, static obstacles, and start-goal queries
    \item Diverse, collision-free demonstration trajectories; disjoint train/validation/test queries
    \item Comparison with the CMA-ES seminar on a common benchmark
\end{itemize}
\subsection{Methods and evaluation}
\begin{itemize}
    \item Simplified Diffuser-style model jointly generating states and actions, with start-goal inpainting
    \item Inpainting-only baseline versus additional obstacle-cost guidance; CMA-ES on an equivalent task
    \item Feasibility (collisions, dynamics, speed bound), plan quality, route diversity, and inference runtime
\end{itemize}
\subsection{Ablations and failure analysis}
\begin{itemize}
    \item Main ablation: contribution of obstacle-cost guidance beyond the learned trajectory prior
    \item Role of training-data coverage; representative collision and dynamics failures
    \item Limits of conclusions drawn from a simple 2D setting
\end{itemize}

\section{Discussion and conclusion}
\begin{itemize}
    \item Comparison of the three papers and the completed 2D experiment
    \item Trade-offs among data requirements, flexibility, feasibility, solution quality, and computational cost
    \item Limitations of learned constraints and iterative sampling; open questions for real-world robot control
\end{itemize}

\bibliographystyle{plain}
\bibliography{Literatur}
\end{document}